% !TeX program = pdflatex
% Technical report: Computer Vision & Visual Localization (EV platform)
% Compile: pdflatex AI_ENGINEER_TECHNICAL_REPORT.tex (run twice for references/hyperref)

\documentclass[12pt,a4paper]{article}

\usepackage[T1]{fontenc}
\usepackage[utf8]{inputenc}
\usepackage[vietnamese,english]{babel}
\selectlanguage{vietnamese}

\usepackage[
  top=2.5cm,
  bottom=2.5cm,
  left=3cm,
  right=2.5cm
]{geometry}

\usepackage{amsmath,amssymb}
\usepackage{graphicx}
\usepackage{booktabs}
\usepackage{xcolor}
\usepackage{float}
\usepackage{caption}
\usepackage{subcaption}
\usepackage{enumitem}
\usepackage{listings}
\usepackage{tikz}
\usetikzlibrary{arrows.meta,positioning,fit,calc}

\usepackage[
  colorlinks=true,
  linkcolor=blue,
  urlcolor=blue,
  citecolor=blue
]{hyperref}

\newcommand{\placeholder}[1]{\colorbox{yellow}{\textbf{[#1]}}}

\lstdefinestyle{pythonstyle}{
  language=Python,
  basicstyle=\ttfamily\small,
  keywordstyle=\color{blue!70!black},
  commentstyle=\color{gray!60},
  stringstyle=\color{teal!60!black},
  showstringspaces=false,
  breaklines=true,
  frame=single,
  tabsize=4,
  columns=fullflexible
}
\lstset{style=pythonstyle}

\title{\textbf{Báo cáo kỹ thuật: Thị giác máy tính \& Định vị thị giác trên nền tảng xe điện}\\[0.5em]
\large \textit{Computer Vision \& Visual Localization --- Engineering implementation}}
\author{\placeholder{TÊN BẠN}}
\date{\today}

\begin{document}
\maketitle

\begin{abstract}
\noindent
Báo cáo mô tả các thành phần đã \textbf{xây dựng và triển khai} trong khuôn khổ hai hướng: (A) phát hiện ổ gà, ước lượng độ sâu/diện tích, phân loại mức độ trên \emph{ONNX Runtime} và CPU; (B) tích hợp ROS\,2 với cơ sở dữ liệu landmark thị giác, \emph{visual odometry}, giám sát toàn vẹn GPS, bộ lọc EKF và các node fusion. Phần cuối phân tích các module gặp hạn chế hoặc chưa hoàn thiện theo hướng kỹ thuật, không dùng bảng chấm điểm hay ngưỡng KPI.
\end{abstract}

\section{Tổng quan hệ thống}

\subsection{Bài toán và hướng tiếp cận}
Hệ thống nhằm hỗ trợ cảm nhận môi trường cho xe điện: (1) phát hiện tổn thương mặt đường (pothole), suy ra độ sâu và diện tích bề mặt ở mức heuristic/metric xấp xỉ; (2) bổ sung định vị khi GPS suy giảm nhờ landmark thị giác, VO và fusion trạng thái. Tiếp cận tổng thể ưu tiên \textbf{triển khai thực tế trên CPU}, dùng \textbf{ONNX} cho inference và ROS\,2 cho phần định vị/fusion.

\subsection{Sơ đồ pipeline tổng thể}
Hình~\ref{fig:pipeline} tóm tắt luồng Part\,A (perception) và Part\,B (localization/fusion) ở mức khối.

\begin{figure}[H]
\centering
\begin{tikzpicture}[
  font=\small,
  box/.style={rectangle, draw, rounded corners, minimum width=2.6cm, minimum height=0.85cm, align=center},
  arr/.style={-{Stealth}, thick}
]
\node[box, fill=blue!8] (in) {Camera / Video\\frame};
\node[box, fill=green!10, right=1.1cm of in] (det) {YOLOv8\\ONNX detect};
\node[box, fill=green!10, right=1.1cm of det] (dm) {Depth map\\(ONNX / gray)};
\node[box, fill=orange!12, right=1.1cm of dm] (geo) {ROI depth\\+ BEV area};
\node[box, fill=orange!12, right=1.1cm of geo] (sev) {Severity\\rules};
\draw[arr] (in) -- (det);
\draw[arr] (det) -- (dm);
\draw[arr] (dm) -- (geo);
\draw[arr] (geo) -- (sev);

\node[box, fill=violet!10, below=1.4cm of det] (img) {Image + GPS\\topics (ROS\,2)};
\node[box, fill=violet!10, right=1.0cm of img] (vo) {VO / lane\\nodes};
\node[box, fill=violet!10, right=1.0cm of vo] (gpsm) {GPS integrity\\state machine};
\node[box, fill=violet!10, right=1.0cm of gpsm] (ekf) {EKF\\fusion};
\draw[arr] (img) -- (vo);
\draw[arr] (vo) -- (gpsm);
\draw[arr] (gpsm) -- (ekf);

\node[draw, dashed, inner sep=6pt, fit=(det)(dm)(geo)(sev), label=above:\textbf{Part A --- Pothole}] {};
\node[draw, dashed, inner sep=6pt, fit=(img)(vo)(gpsm)(ekf), label=below:\textbf{Part B --- GPS + Visual}] {};
\end{tikzpicture}
\caption{Pipeline tổng thể: perception pothole (trên) và khối ROS\,2 localization/fusion (dưới).}
\label{fig:pipeline}
\end{figure}

\subsection{Môi trường chạy}
\begin{itemize}[leftmargin=*]
  \item \textbf{OS:} Windows (phát triển); script Python và workspace ROS\,2 được tổ chức dưới thư mục dự án \texttt{gps\_system} (trước đây \texttt{ros2\_ws}) và \texttt{pothole}.
  \item \textbf{Python:} \placeholder{Phiên bản Python (vd.\ 3.10)}; inference pothole dùng \texttt{onnxruntime} khi bật depth ONNX; detector ONNX có thể qua Ultralytics decode.
  \item \textbf{Compute:} ưu tiên \textbf{CPU-only} cho pipeline pothole; GPU tùy chọn khi có CUDA cho Ultralytics/ORT.
  \item \textbf{ROS\,2:} build bằng \texttt{colcon} cho package \texttt{gps\_visual\_b} (nodes, launch, params).
\end{itemize}

%---------------------------------------------------------------------
\section{Phần A: Phát hiện và đo độ sâu ổ gà}

\subsection{Object Detection}
\subsubsection*{Kiến trúc model và dữ liệu}
Detector dựa trên \textbf{YOLOv8} fine-tune trên dataset nội bộ với \textbf{3 lớp} (\texttt{minor}, \texttt{moderate}, \texttt{severe}), cấu hình trong \texttt{dataset.yaml} (tách \texttt{train/val/test}). Quy trình train/export nằm trong repo (\texttt{train\_yolo}, \texttt{export\_onnx}).

\subsubsection*{ONNX export và inference}
Trọng số được export sang \texttt{.onnx} để chạy inference thống nhất trên ONNX Runtime hoặc qua lớp bọ Ultralytics \texttt{YOLO(best.onnx)} để decode chuẩn đầu ra YOLOv8, tránh lệch định dạng tensor thủ công.

\subsubsection*{Kết quả đạt được (mô tả, không dùng bảng điểm)}
Trên tập test detector, một lần đánh giá Ultralytics đã ghi trong báo cáo nội bộ các giá trị sau (xuất JSON/PNG, không dùng làm bảng chấm điểm ở đây): \emph{mAP@0.5} khoảng $0.72$, \emph{mAP@0.5:0.95} khoảng $0.43$, precision khoảng $0.78$ và recall khoảng $0.66$. Đường cong PR và confusion matrix được lưu kèm (\texttt{results.png}, \texttt{confusion\_matrix.png}). Tốc độ xử lý end-to-end phụ thuộc cấu hình (CPU vs GPU); ví dụ log pipeline video có thể đạt cỡ hàng chục FPS trung bình khi cấu hình nhẹ, giảm rõ khi bật depth ONNX nặng.

Phân loại \textbf{severity} không phải head phân lớp thứ tư của YOLO mà là \textbf{rule-based} sau đo lường: ngưỡng trong \texttt{severity\_rules.yaml} kết hợp \texttt{depth\_m} và \texttt{area\_m2}.

\begin{lstlisting}[caption={Luồng xử lý cốt lõi theo frame (rút gọn từ \texttt{PotholePipeline})}]
# depth_map: tu Depth ONNX hoac fallback tu grayscale
depth_map = self._estimate_depth_map(frame)
detections = self.detector.infer(frame)
for det in detections:
    area_m2 = estimate_area_m2(frame, det.bbox_xyxy, mpp_bev)
    depth_m = estimate_depth_m(depth_map, det.bbox_xyxy, camera_cfg)
    severity = classify_severity(depth_m, area_m2, severity_cfg)
\end{lstlisting}

\subsection{Depth và Area Estimation}
\subsubsection*{Hướng tiếp cận}
Chọn \textbf{monocular depth} (Depth Anything V2 dạng ONNX khi có) vì phù hợp một camera phía trước; khi thiếu model nặng, dùng \textbf{relative depth} cực nhẹ từ biến thể grayscale để giữ real-time trên CPU.

\subsubsection*{Công thức và thuật toán}
\paragraph{Depth (ROI).}
Cho depth map tương đối $D$ và bbox $(x_1,y_1,x_2,y_2)$, lấy median trong ROI và vùng ngữ cảnh mặt đường phía trên cạnh trên bbox; độ lõm tương đối $\max(0,\,\mathrm{median}_{\mathrm{roi}}-\mathrm{median}_{\mathrm{road}})$ được nhân với hệ số gắn với \texttt{camera\_height\_m} để ra \texttt{depth\_m} heuristic (xem \texttt{estimate\_depth\_m} trong \texttt{geometry.py}).

\paragraph{Diện tích (BEV xấp xỉ).}
Homography gần đúng từ bốn điểm mặt đường trong ảnh xuống mặt phẳng BEV, \texttt{warpPerspective}, mask Otsu trên ROI BEV, đếm pixel nhân với $(\texttt{meters\_per\_pixel\_bev})^2$.

\subsubsection*{Calibration}
Tham số trong \texttt{camera\_config.yaml}: \texttt{camera\_height\_m}, \texttt{meters\_per\_pixel\_bev} điều chỉnh scale metric; homography hiện là \textbf{template geometry} (chưa full extrinsic calibration ngoài hiện trường).

\subsubsection*{Kết quả định tính}
Có script \texttt{evaluate\_test\_estimation.py} so bbox GT (CSV) với ước lượng lại trên ảnh test sau scale không gian; output gồm \texttt{per\_detection\_errors.csv} và \texttt{summary\_metrics.json}. Sai số phụ thuộc mạnh vào đúng không gian ảnh tham chiếu của CSV và chất lượng depth tương đối.

\subsection{ROI Fusion (Bbox $\cap$ Depth)}
Với mỗi detection, bbox được clip trong giới hạn ảnh; \textbf{cùng một depth map} (tái sử dụng theo chu kỳ \texttt{depth\_interval} trong pipeline) được cắt ROI đúng tọa độ integer để tính median, tránh tính lại depth map mỗi detection. Đây là thiết kế \textbf{đồng bộ trong một luồng} trên CPU; tùy chọn giảm tần suất depth để cân bằng FPS.

%---------------------------------------------------------------------
\section{Phần B: GPS và dự phòng thị giác}

\subsection{Visual Landmark Database}
Package \texttt{gps\_visual\_b} chứa \texttt{landmark\_database.py} và node \texttt{landmark\_db\_node}: lưu trữ landmark dạng keyframe kèm descriptor phục vụ ghép đồng nhất đoạn đường. Descriptor có thể mở rộng (VPR encoder trong repo); cấu trúc tách module để bổ sung dữ liệu đặc trưng theo tuyến.

\subsection{Visual Odometry / Lane}
\texttt{visual\_odometry.py} và \texttt{visual\_odometry\_node.py} cung cấp khung VO; \texttt{lane\_bev.py} / \texttt{lane\_detection\_node.py} hỗ trợ làn đường trong BEV. Scale drift được xử lý ở mức thiết kế bằng coupling với GPS tốt và EKF (xem dưới), chứ chưa thay thế SLAM metric đầy đủ.

\subsection{GPS Integrity Monitor và Handover}
Module \texttt{gps\_integrity.py} cài đặt máy trạng thái \texttt{GPS\_GOOD} $\rightarrow$ \texttt{GPS\_DEGRADED} $\rightarrow$ \texttt{GPS\_LOST} dựa trên cờ fix, HDOP, số vệ tinh (và mở rộng SNR nếu có topic phụ). Có cơ chế \textbf{latched fix} WGS84 tin cậy cao để phục vụ handover khi mất tín hiệu ngắn.

\subsection{Sensor Fusion Engine}
\texttt{ekf\_fusion.py} triển khai bộ lọc \textbf{EKF} cho vector trạng thái liên quan tới pose/velocity; \texttt{sensor\_fusion\_node.py} tích hợp ROS\,2. Cấu trúc node trong \texttt{launch/system\_b.launch.py} và \texttt{config/params.yaml} điều phối các node monitor, lane, VO, fusion.

%---------------------------------------------------------------------
\section{Failure Analysis}

\subsection*{FM-01: False negative ổ gà nhỏ / nông hoặc phản xạ mạnh}
\textbf{Vấn đề gặp phải:} ổ gà thấy được nhưng không có bbox; tăng khi mặt đường ướt hoặc tương phản thấp.\\
\textbf{Nguyên nhân:} domain shift so với train; ngưỡng confidence; object nhỏ sau resize.\\
\textbf{Đã thử:} điều chỉnh ngưỡng, augmentation trong tài liệu đề xuất.\\
\textbf{Hướng khắc phục:} thu thập dữ liệu điều kiện mưa/đêm; hard-negative mining; multi-scale test-time.

\subsection*{FM-02: False positive (vệt vá, vũng sáng)}
\textbf{Vấn đề gặp phải:} bbox trên vùng không phải pothole.\\
\textbf{Nguyên nhân:} texture tương tự; thiếu post-filter hình thái.\\
\textbf{Đã thử:} rule severity không giải quyết FP thuần detection.\\
\textbf{Hướng khắc phục:} filter theo aspect ratio, tracking temporal, secondary classifier.

\subsection*{FM-03: FPS overlay vs độ mượt cảm nhận}
\textbf{Vấn đề gặp phải:} FPS cao nhưng video vẫn giật.\\
\textbf{Nguyên nhân:} đo cũ không phản ánh wall-clock end-to-end.\\
\textbf{Đã thử:} chuyển đo e2e, FPS trượt, hiển thị \texttt{PROC ms}.\\
\textbf{Hướng khắc phục:} tối ưu I/O và giảm độ phân giải hiển thị.

\subsection*{FM-04: Sụt hiệu năng khi bật depth ONNX lớn}
\textbf{Vấn đề gặp phải:} FPS giảm mạnh trên CPU.\\
\textbf{Nguyên nhân:} backbone depth lớn; pipeline tuần tự.\\
\textbf{Đã thử:} \texttt{depth\_interval}, giảm \texttt{depth\_input\_size}.\\
\textbf{Hướng khắc phục:} INT8 quantize; thread tách depth; model nhỏ hơn.

\subsection*{FM-05: ONNX dynamic shape depth}
\textbf{Vấn đề gặp phải:} lỗi symbolic height/width hoặc broadcast.\\
\textbf{Nguyên nhân:} export dynamic không tương thích mọi size.\\
\textbf{Đã thử:} fallback input cố định (518) trong \texttt{DepthOnnxEstimator}.\\
\textbf{Hướng khắc phục:} export cố định kích thước hoặc simplify graph.

\subsection*{FM-06: Sai số depth/area do mô hình hình học xấp xỉ}
\textbf{Vấn đề gặp phải:} độ sâu/diện tích dao động theo góc nhìn.\\
\textbf{Nguyên nhân:} monocular relative; homography template chưa calibrate ngoại tại.\\
\textbf{Đã thử:} rule severity + logging.\\
\textbf{Hướng khắc phục:} extrinsic calibration; stereo; scale từ IMU/wheel.

\subsection*{Phần B: tích hợp ROS\,2 trên Windows}
\textbf{Vấn đề gặp phải:} log colcon từng báo lỗi symlink/privilege khi build.\\
\textbf{Nguyên nhân:} khác biệt filesystem Windows vs Linux trong workflow ROS.\\
\textbf{Đã thử:} build lại, điều chỉnh workspace path (\texttt{gps\_system}).\\
\textbf{Hướng khắc phục:} WSL2 hoặc CI Linux cho build sản phẩm.

%---------------------------------------------------------------------
\section*{Phụ lục: Phiên bản phần mềm (tham chiếu)}
\begin{table}[H]
\centering
\begin{tabular}{@{}ll@{}}
\toprule
Thành phần & Ghi chú \\
\midrule
ONNX Runtime & Inference depth/detector (tùy cấu hình) \\
Ultralytics & Decode YOLOv8 ONNX (detector) \\
OpenCV & Geometry, BEV, visualization \\
ROS\,2 (Humble / setup tùy máy) & Nodes Part\,B \\
\bottomrule
\end{tabular}
\caption{Bảng tham chiếu công cụ (không phải bảng chấm điểm hay KPI).}
\end{table}

\vfill
\noindent\textit{End of document.}

\end{document}
